\chapter{Enveloping bars and planar path sequences}
\label{ch:enveloping-audit}

The dimension of a bar chain space and the dimension of its homology are different invariants. For an enveloping algebra, the Chevalley comparison gives a Lie-theoretic model for the homology. This makes it possible to compare Motzkin and Riordan path counts with the ordered bar on a specified carrier.

Bar length and Chevalley degree are homological degrees: $H_n$ denotes degree-$n$ homology. The cohomology of the dual Chevalley complex is written $H^n$. For finite-dimensional weight pieces their dimensions agree.

\section{Chevalley comparison}

\begin{theorem}[Bar--Chevalley comparison]
\label{thm:bar-chevalley-spine}
Let $\g$ be a Lie algebra object in $\FactD(X)$ for the pointwise product,
$K$-flat, positively graded and complete, with every total weight admitting
finitely many decompositions into positive weights, and suppose internal PBW
holds for $U(\g)$. There is a natural quasi-isomorphism of conilpotent
coalgebras in $\FactD(X)$
\[
 \BarPerp\bigl(U(\g)\bigr)\ \xrightarrow{\ \sim\ }\ C_*(\g),
\]
compatible with the factorization coefficient and the PBW filtration.
\end{theorem}
\begin{proof}
Filter $U(\g)$ by PBW degree and both complexes by total word length. The
associated graded enveloping algebra is $\Sym(\g)$, and the bar of a symmetric
algebra is the suspended symmetric coalgebra on $\g$, which is the associated
graded of Chevalley chains. The comparison is an isomorphism on the first page;
finite decomposition gives a bounded exhaustive spectral sequence in each
weight, and completeness promotes it to a quasi-isomorphism. Every step uses
only internal tensor products and the bracket, so the map stays in $\FactD(X)$.
\end{proof}

\begin{calculation}[Finite-dimensional simple case]
\label{calc:simple-spine}
For a simple Lie algebra of rank $r$ with exponents $m_1,\dots,m_r$, the
Chevalley cohomology is the exterior algebra on primitive classes of degrees
$2m_i+1$. The corresponding homology Poincar\'e polynomials are
\[
 \BarPerp U(\mathfrak{sl}_2):\ 1+t^3,
 \qquad
 \BarPerp U(\mathfrak{sl}_3):\ 1+t^3+t^5+t^8,
 \qquad
 \BarPerp U(\mathfrak g_2):\ 1+t^3+t^{11}+t^{14}.
\]
The corresponding Betti sequences for $\mathfrak{sl}_2$ and $\mathfrak{sl}_3$ are
$1,0,0,1$ and $1,0,0,1,0,1,0,0,1$.
For a finite-dimensional Lie algebra, the Chevalley chain dimensions are
$\binom{\dim\g}{n}$.

For an ordinary augmented algebra $A$ over $\mathbb C$, let
$I=\ker(A\to\mathbb C)$.  Its normalized two-sided bar with trivial
coefficients has homological chain spaces
\[
 B_0(A)=\mathbb C,\qquad B_n(A)=I^{\otimes n}\quad(n\ge1).
\]
Here $n$ is bar length, independently of any internal weight or PBW degree.
If $I$ has finite dimension $d$, then $\dim B_n(A)=d^n$; in particular,
an augmentation ideal of dimension three gives the chain counts $3^n$.
In a cohomological convention that places $B_n$ in degree $-n$, its homology
is denoted $H^{-n}$ instead of $H_n$.

For $A=U(\mathfrak{sl}_2)$, PBW gives the basis
$e^a h^b f^c$, with $a,b,c\ge0$, where
$[h,e]=2e$, $[h,f]=-2f$, and $[e,f]=h$.
The augmentation ideal has all these basis elements except $1$.
Consequently $B_n(A)$ is countably infinite-dimensional for every $n\ge1$.
Its total PBW filtration has finite-dimensional associated graded pieces,
with generating function
\[
 \sum_{d\ge0}\dim\operatorname{gr}^{\mathrm{PBW}}_d B_n(A)\,z^d
 =\bigl((1-z)^{-3}-1\bigr)^n.
\]
In particular, $3^n$ is the dimension in the least possible total PBW degree
$d=n$. These least-degree dimensions are $3,9,27,81,243$ for
$n=1,2,3,4,5$; the full bar chain spaces are infinite-dimensional. The quotient of
indecomposables is zero: the perfectness of $\mathfrak{sl}_2$ implies $I/I^2=0$.

The Chevalley chain model is different: its degree-$n$ space is
$\Lambda^n\mathfrak{sl}_2$, of dimensions $1,3,3,1$ in degrees $0,1,2,3$.
With the convention $d(x\wedge y)=[x,y]$, the map
$d:\Lambda^2\mathfrak{sl}_2\to\mathfrak{sl}_2$ sends
\[
 e\wedge h\longmapsto-2e,\qquad
 e\wedge f\longmapsto h,\qquad
 h\wedge f\longmapsto-2f,
\]
and is an isomorphism.  The next differential vanishes, since
\[
 d(e\wedge h\wedge f)
 =-2e\wedge f-h\wedge h-2f\wedge e=0.
\]
Together with $d:\mathfrak{sl}_2\to\mathbb C$ equal to zero, this gives
Chevalley homology dimensions $1,0,0,1$.  The ordinary bar--Chevalley
comparison identifies these homology groups with those of $U(\mathfrak{sl}_2)$;
it does not identify the two chain spaces.

\end{calculation}

\section{The Virasoro row}

The transverse presentation of the Virasoro family is the enveloping algebra of
the positive Witt algebra
\[
 L_1=\bigoplus_{n\ge1}\mathbb C\,e_n,
 \qquad
 [e_i,e_j]=(j-i)e_{i+j},
\]
to which Theorem~\ref{thm:bar-chevalley-spine} applies: $L_1$ is positively
graded with one-dimensional weight spaces, so every weight has finitely many
decompositions.

\begin{theorem}[Pentagonal rigidity]
\label{thm:pentagonal-spine}
The bar homology of $U(L_1)$ satisfies
\[
 \dim H_n\bigl(\BarPerp U(L_1)\bigr)=2
 \qquad\text{for every }n\ge1,
\]
one dimension in each of the weights $(3n^2-n)/2$ and $(3n^2+n)/2$. Its graded
Euler character is
\begin{equation}
\label{eq:witt-euler-spine}
 \chi_q\bigl(\BarPerp U(L_1)\bigr)=\prod_{n\ge1}\bigl(1-q^n\bigr),
\end{equation}
and \eqref{eq:witt-euler-spine} is accounted for degree by degree with no
residual freedom.
\end{theorem}
\begin{proof}
By Theorem~\ref{thm:bar-chevalley-spine} the bar of $U(L_1)$ is quasi-isomorphic to $C_*(L_1)$.
Each weight subcomplex is finite-dimensional, so its homology has the same
dimensions as the weightwise dual cohomology. Goncharova's calculation gives
two dimensions in each degree $n\ge1$, in the weights $(3n^2\mp n)/2$. Identity \eqref{eq:witt-euler-spine} is the
denominator product applied to $L_1$, where every root space is one dimensional
and even, so $\smult(L_1)_n=1$. Euler's pentagonal number theorem expands the
product as $\sum_{m\in\Z}(-1)^mq^{m(3m-1)/2}$, whose support is exactly
$\{(3n^2\pm n)/2\}$ with coefficient $(-1)^n$ at each of $(3n^2-n)/2$ and
$(3n^2+n)/2$. Comparing with $\sum_n(-1)^n\dim H_{n,w}$, each weight of the Euler
character is saturated by a single class in the degree Goncharova predicts, so
no further classes and no cancellation are available.
\end{proof}

For $n=1,2,3,4$, the nonzero homology groups occupy the weight pairs
$(1,2)$, $(5,7)$, $(12,15)$, $(22,26)$, one dimension in each.
The corresponding nonzero coefficients of $\prod_{n\ge1}(1-q^n)$ occur at
$0,1,2,5,7,12,15,22,26$ with signs $+,-,-,+,+,-,-,+,+$.

\section{Planar path sequences}

A Motzkin path is a lattice path from $(0,0)$ to $(n,0)$, staying at
nonnegative height, with steps $(1,1)$, $(1,0)$, and $(1,-1)$.
Let $M(n)$ count all such paths of length $n$, and let $R(n)$ count those with
no horizontal step at height zero.  Both counts include the empty path:
$M(0)=R(0)=1$, while $M(1)=1$ and $R(1)=0$.
For $n\ge1$ consider the two sequences
\[
 a_n=M(n+1)-M(n)=1,2,5,12,30,76,\ldots,
 \qquad
 b_n=R(n+3)=3,6,15,36,91,\ldots.
\]

\begin{proposition}[Motzkin and Riordan identities]
\label{prop:one-sequence-spine}
For every $n\ge0$,
\[
 M(n)=R(n)+R(n+1),\qquad
 M(n+1)-M(n)=R(n+2)-R(n).
\]
The ordinary generating functions of the two path counts, and hence those
of $a_n$ and $b_n$, belong to
$\mathbb Q(z,\sqrt{1-2z-3z^2})$.
\end{proposition}
\begin{proof}
Work in the formal power-series ring $\mathbb Q[[z]]$, and write
$M(z)=\sum_{n\ge0}M(n)z^n$ and $R(z)=\sum_{n\ge0}R(n)z^n$.
A nonempty Motzkin path either starts with a horizontal step followed by
an arbitrary Motzkin path, or has a unique first-return decomposition
$U P D Q$, where $U$ and $D$ are the up and down steps,
$P$ and $Q$ are arbitrary Motzkin paths, and $P$ is drawn
one unit above the horizontal axis.  Thus
\[
 M=1+zM+z^2M^2.
\]
A nonempty path counted by $R$ has no initial horizontal step.  Its first
return similarly has the form $U P D Q$, with $P$ arbitrary and $Q$ counted
by $R$.  Consequently
\[
 R=1+z^2MR,\qquad R=(1-z^2M)^{-1}.
\]
The inverses here exist because their denominators have constant term one.
Using $z^2M^2=M-1-zM$, compute
\[
 (1+zM)(1-z^2M)
 =1+zM-z^2M-z^3M^2=1+z.
\]
It follows that $(1+z)R=1+zM$.  Comparing coefficients of $z^{n+1}$ gives
$R(n+1)+R(n)=M(n)$ for all $n\ge0$.  Subtracting consecutive identities
gives the difference formula.

Solving the quadratic equation for the formal series with constant term
one gives
\[
 M(z)=\frac{1-z-\sqrt{1-2z-3z^2}}{2z^2},\qquad
 R(z)=\frac{1+zM(z)}{1+z},
\]
where the square root has constant term one.  The apparent pole in the
formula for $M$ is removable as a formal series.  Finite index shifts and
differences preserve the indicated quadratic function field, proving the
last assertion.  These identities are consequences solely of the path
decompositions.
\end{proof}

\begin{remark}[Low-degree homology]
Theorem~\ref{thm:pentagonal-spine} gives homology dimensions $2,2,2$ in degrees $1,2,3$, whereas $a_1,a_2,a_3$ are $1,2,5$. These triples agree only in degree two. The degree-one discrepancy already excludes equality of the full homology dimension sequences.
\end{remark}

\begin{theorem}[Enveloping bars and planar path counts]
\label{thm:carrier-spine}
Let $B_\bullet(A)$ be the ordinary normalized bar over $\mathbb C$, with
trivial coefficients and homological degree given by bar length.  With
$a_n=M(n+1)-M(n)$ and $b_n=R(n+3)$ indexed by $n\ge1$ as above,
\[
 \dim H_1 B_\bullet(U(\mathfrak{sl}_2))=0\ne b_1=3,
 \qquad
 \dim H_1 B_\bullet(U(L_1))=2\ne a_1=1,
\]
where $L_1=\bigoplus_{m\ge1}\mathbb C e_m$ and
$[e_i,e_j]=(j-i)e_{i+j}$.
Hence neither sequence is the bar homology dimension sequence of the
corresponding enveloping algebra.  Neither is its full bar chain dimension
sequence: in both cases $B_n$ is countably infinite-dimensional for every
$n\ge1$.
\end{theorem}
\begin{proof}
For any augmented algebra $A$ with augmentation ideal $I$, the bar
differential out of $B_1=I$ is zero and the image from $B_2=I\otimes I$ is
$I^2$.  Therefore
\[
 H_1B_\bullet(A)=I/I^2.
\]
For an ordinary Lie algebra $\mathfrak g$, the inclusion of generators
induces an isomorphism
\[
 \mathfrak g/[\mathfrak g,\mathfrak g]\ \xrightarrow{\ \sim\ }
 I/I^2,\qquad I=\ker(U(\mathfrak g)\to\mathbb C).
\]
Indeed the generators span $I/I^2$, and
$[x,y]=xy-yx$ vanishes there.  For injectivity, equip
$\mathbb C\oplus\mathfrak g/[\mathfrak g,\mathfrak g]$ with the square-zero
product on its second summand.  The quotient map from $\mathfrak g$ to that
summand is a Lie map into the commutator Lie algebra of this associative
algebra.  The universal property of $U(\mathfrak g)$ extends it to an
augmented algebra map, which kills $I^2$ and provides an inverse on the
displayed quotient.

The brackets $[h,e]=2e$, $[h,f]=-2f$, and $[e,f]=h$ show that
$\mathfrak{sl}_2$ equals its commutator algebra.  For $L_1$, every nonzero
bracket has weight at least three.  Conversely, for $m\ge3$,
\[
 e_m=\frac{1}{m-2}[e_1,e_{m-1}].
\]
Thus $[L_1,L_1]=\bigoplus_{m\ge3}\mathbb C e_m$, and its abelianization
has basis the classes of $e_1,e_2$, in weights one and two.  This proves
both degree-one assertions without using any higher-degree calculation.

PBW gives a countably infinite basis of the augmentation ideal of either
enveloping algebra.  In the $L_1$ case the basis consists of nonempty
ordered monomials in the countable set $e_1,e_2,\ldots$, each monomial
having finite length.  Finite tensor powers of either augmentation ideal
are therefore countably infinite-dimensional.  The path counts are finite
in every degree, so they cannot count these full chain spaces.
\end{proof}

In particular, a complex whose homology dimensions are $b_n$ cannot be
quasi-isomorphic to $B_\bullet(U(\mathfrak{sl}_2))$, and one whose homology
dimensions are $a_n$ cannot be quasi-isomorphic to $B_\bullet(U(L_1))$,
with the stated degree convention.  A discrepancy of chain dimensions
alone would not imply this conclusion; the degree-one homology discrepancy
does.

The weight grading gives a different, finite comparison.  An ordered PBW
monomial of weight $w$ specifies a partition of $w$, so for $n\ge1$,
\[
 \sum_{w\ge0}\dim B_n(U(L_1))_w q^w
 =\left(\prod_{m\ge1}(1-q^m)^{-1}-1\right)^n.
\]
Each coefficient is finite.  Thus a weightwise dimension, a dimension
after a specified weight cutoff, and an ungraded dimension are distinct
invariants.  In a weight completion, dimensions must likewise be interpreted
weightwise; the countable algebraic dimensions above refer to the ordinary
direct-sum enveloping algebras.

\begin{remark}[Chiral coefficients and transverse multiplication]
Independently of these scalar calculations, a Virasoro chiral coefficient does not determine the transverse bar: a fixed coefficient may be combined with free, square-zero, enveloping, or quantum-plane transverse algebras whose bars differ, while the coefficient and the central charge are unchanged. A chain-level comparison with an ordered bar therefore also requires a specified transverse multiplication and a quasi-isomorphism from the proposed complex to that bar.
\end{remark}

\chapter{The obstruction constitution}
\label{ch:obstructions}

Each statement below names what is \emph{not} determined, and therefore what
must be supplied. Without them the preceding theorems admit false corollaries.

\begin{proposition}[Where a mixed distributive law can live]
\label{prop:lambda-sector-spine}
The sector on which a mixed distributive law between $\tensor^!$ and $\starCh$
is geometrically defined is the common-decomposition sector, the image of the
aligned idempotent $e_{\mathrm{al}}$. Off that sector the two products admit no
interchange, only the aligned retraction, so no distributive law can be
written.
\end{proposition}
\begin{proof}
A distributive law requires a natural transformation interchanging the two
products in a fixed direction, subject to the coherence conditions of each. The
canonical comparison in the additive model is the aligned inclusion $\iota$ with
retraction $\pi$, and $\pi\iota=\id$ while $\iota\pi=e_{\mathrm{al}}\ne\id$ in
general. A natural transformation subject to both coherence conditions
therefore exists only where $e_{\mathrm{al}}$ acts as the identity.
\end{proof}

This locates a mixed law without constructing one. No example in the present
corpus exhibits one; Zamolodchikov--Faddeev algebras are the natural candidate,
since their spectral variable belongs to the chiral direction while their word
belongs to the transverse direction.

\begin{theorem}[No quantum group from dimension loss]
\label{nogo:dimension-loss-spine}
The dimension of the kernel of $V^{\otimes n}\to\Sym^nV$ determines no
$R$-matrix, associator, or quantum group. That kernel carries no
multiplication, coproduct, monodromy, or coherence equation.
\end{theorem}
\begin{proof}
The kernel is a representation of $\Sigma_n$ and nothing more. Two algebras
with the same underlying graded vector space, hence the same kernel dimensions
in every arity, may have different multiplications and inequivalent
representation categories; a free algebra and a quadratic quotient with the
same Hilbert series is an example. A quantum-group structure is a coproduct
together with coherence data, none of which is a function of a dimension.
\end{proof}

\begin{theorem}[No modular object from topology alone]
The identity $\partial^2=0$ on $\overline{\mathcal M}_{g,n}$ assigns no
coefficient operations to boundary strata. A modular master element exists only
after cyclic contractions and clutching-compatible coefficient maps have been
supplied.
\end{theorem}

\begin{theorem}[No bulk from a boundary centre]
The derived centre is terminal among algebraic central actions on the boundary
algebra. It does not detect a decoupled bulk sector, so a physical bulk is not
determined by the boundary chiral algebra unless a bulk-to-centre comparison is
supplied and proved to be an equivalence.
\end{theorem}
\begin{proof}
Terminality is the universal property of the centre. It is a statement about
the boundary category and is insensitive to any sector acting trivially on it,
so a bulk with trivial boundary action is invisible.
\end{proof}

\begin{theorem}[No classification by central charge]
Central charge determines the Virasoro anomaly coefficient. It does not
determine the chiral operation, the representation category, the braid
monodromy, or the higher-genus conformal blocks. Distinct holomorphic theories
at equal central charge exist.
\end{theorem}

\begin{theorem}[No identification of anomalies of different types]
\label{nogo:anomalies-spine}
A current-algebra level, a Virasoro central charge, a determinant-line
curvature, a BRST ghost-number anomaly, and a Borcherds-product weight are
invariants in different theories, living in different deformation complexes.
They may be related by a theorem in a particular family. They are not one
universal scalar, and no identity among them follows from numerical agreement.
\end{theorem}
\begin{proof}
Each is a class in a distinct complex: an Atiyah-algebra extension class, a
Virasoro cocycle, the curvature of a determinant line, a class in $H^1$ or
$H^2$ of a BRST convolution complex, and a weight in the character lattice of
an automorphic form. A map between two of these complexes is additional data;
in its absence, equality of numbers is a coincidence of values, not an
identification of classes.
\end{proof}

\begin{theorem}[No unnamed comparison of bars]
An asserted equality $\BarPerp(A)=\BarAssCh(A)$ has no meaning without a
functor carrying the pointwise multiplication problem to the chiral-convolution
multiplication problem. The two have different underlying gradings, counting
interval subdivisions and collision trees respectively. Agreement of Euler
characteristics is not a comparison map.
\end{theorem}

\begin{theorem}[A scalar sequence is not homology data]
\label{nogo:scalar-spine}
A sequence of integers becomes bar homology only after a chain complex, a
differential, a grading convention, and a rank computation or quasi-isomorphism
are specified. Matching initial values, discriminants, or growth rates
constructs no such map.
\end{theorem}
\begin{proof}
Chain-space dimensions and homology dimensions differ by the ranks of the
differentials, which are not functions of the chain dimensions.
Calculation~\ref{calc:simple-spine} exhibits an algebra whose bar chain
dimensions are $3,9,27,81,243$ and whose bar homology dimensions are $1,0,0,1$;
Proposition~\ref{prop:one-sequence-spine} exhibits two sequences with the same
discriminant and growth rate that are transforms of one another and of neither
homology.
\end{proof}

Theorem~\ref{nogo:anomalies-spine} together with the central-charge statement
retracts any programme treating a tuple of scalars drawn from these sources as
a single classifying invariant, and any trace identity asserted among them
without a named comparison morphism.
